\section{Almost divisibility with a relative error}
\label{ext:sec:divisibility}

For a given class $x$ and integer $N\ge1$, we seek a single class $u$
with $\varphi(x)/(N+1)<\varphi(u)<\varphi(x)/N$ whenever
$0<\varphi(x)<\infty$. We choose $s$ between $1/(N+1)$ and $1/N$
and approximate $s\varphi(x)$ from below. The error is measured by one
fixed cut, so the same construction also controls functionals taking the
value infinity on $x$.

\begin{theorem}
\label{ext:thm:divisibility}
Let $A$ be a simple exact non-elementary $C^*$-algebra with strict
comparison as in Equation~\eqref{ext:eq:comparison}.
For every $x\in\Cu(A)$ and integer $N\ge1$, there exists
$u\in\Cu(A)$ such that $Nu\le x\le(N+1)u$.
\end{theorem}

\begin{proof}
The assertion is immediate for $x=0$. Otherwise write $x=[a]$,
where $a\in B_+$ has norm one, and put $e=(a-1/2)_+$.
If $\Fin(A)$ is empty, all functionals are zero or infinite on
every nonzero class. Equation~\eqref{ext:eq:comparison} then gives
$Nx\le x$, while $x\le(N+1)x$ is automatic. Take $u=x$.
We henceforth assume $\Fin(A)\ne\varnothing$.

\medskip\noindent
\textit{Choose supports with summable relative rank.}
Choose a real number $s$ and integers $M_j\ge1$, $j\ge1$, such that
\begin{equation}
\label{ext:eq:budget}
 \frac1{N+1}<s<\frac1N,\qquad
 E:=\sum_{j=1}^{\infty}\frac1{M_j}<s-\frac1{N+1}.
\end{equation}
Lemma~\ref{ext:lem:small} supplies nonzero $h_j\in B_+$ with
$M_j[h_j]\le[e]$.

To recover the full rank of $a$ from functions vanishing near zero,
choose increasing cutoffs $F_j$ converging to one on $(0,1]$ and
transfer their positive increments. Concretely, set
\[
 F_0=0,\qquad F_j(t)=\min\{1,\max\{0,2^jt-1\}\},\qquad
 r_j=s(F_j-F_{j-1})\quad(j\ge1).
\]
The functions $F_j$ increase, and $r_j\ge0$ vanishes on
$[0,2^{-j}]$. Put $p_j=r_j(a)$.
Lemma~\ref{ext:lem:transfer} gives $c_j\in\her(h_j)_+$ satisfying
\begin{equation}
\label{ext:eq:transferred}
 \tau(c_j)=\tau(p_j)\quad\text{for all extended traces }\tau,
 \qquad M_j[c_j]\le[e].
\end{equation}

\medskip\noindent
\textit{Round and sum inside the Cuntz semigroup.}
Choose integers $L_j>\|c_j\|$ and put
\begin{equation}
\label{ext:eq:u}
 w_j=\sum_{i=1}^{L_j}[(c_j-i)_+],\qquad
 u=\sup_{k\ge1}\sum_{j=1}^k w_j.
\end{equation}
For completeness, this supremum is represented by a positive
element of $B$: choose representatives of $w_j$, rescale them
to norm at most $2^{-j}$, and put them in mutually orthogonal
stable blocks. Their sum converges in norm. Each finite partial
sum is Cuntz below it. Conversely, comparison of all partial sums
with one target passes to the norm limit, proving the least-upper-bound
property. Therefore
\begin{equation}
\label{ext:eq:u-values}
 \varphi(u)=\sum_{j=1}^{\infty}\varphi(w_j)
 \quad(\varphi\in\F(\Cu(A))).
\end{equation}

Fix $\varphi\in\Fin(A)$ and choose its representing trace
using Equation~\eqref{ext:eq:trace-representation}.
Equations~\eqref{ext:eq:rounding} and~\eqref{ext:eq:transferred} give
\begin{equation}
\label{ext:eq:piece-error}
 \varphi(w_j)\le\tau(p_j)
 \le\varphi(w_j)+\frac{\varphi([e])}{M_j}.
\end{equation}
The finite value $\varphi([e])$ is the only trace-dependent
constant in the error estimate.

\medskip\noindent
\textit{Recover the rank of $a$.}
Telescoping gives
\[
 \sum_{j=1}^k\tau(p_j)=s\tau(F_k(a)).
\]
We claim that
\begin{equation}
\label{ext:eq:recover}
 \tau(F_k(a))\uparrow d_\tau(a)=\varphi(x).
\end{equation}
The values increase and are at most $d_\tau(a)$ by
Equation~\eqref{eq:pre:trace-rank} and hereditary containment.
Given $0<\epsilon<1$, for sufficiently large $k$ we have
$F_k(t)=1$ on $[\epsilon,1]$. Thus, simultaneously for all
integers $m\ge1$,
\[
 (a-\epsilon)_+^{1/m}\le F_k(a).
\]
Taking traces and then letting $m\to\infty$ gives
$d_\tau((a-\epsilon)_+)\le\tau(F_k(a))$.
The cut-supremum formula in Equation~\eqref{ext:eq:cu-facts} now proves
Equation~\eqref{ext:eq:recover}.

Sum Equation~\eqref{ext:eq:piece-error}, then use
Equations~\eqref{ext:eq:u-values} and~\eqref{ext:eq:recover}. In extended
nonnegative arithmetic, for every $\varphi\in\Fin(A)$,
\begin{equation}
\label{ext:eq:total-error}
 \varphi(u)\le s\varphi(x)
 \le\varphi(u)+E\varphi([e]).
\end{equation}
If $\varphi(x)<\infty$, Equation~\eqref{ext:eq:controller} implies
\[
 (s-E)\varphi(x)\le\varphi(u)\le s\varphi(x).
\]
Since $\varphi(x)>0$, Equation~\eqref{ext:eq:budget} yields
\begin{equation}
\label{ext:eq:strict-interval}
 \frac{\varphi(x)}{N+1}<\varphi(u)<\frac{\varphi(x)}N.
\end{equation}
If instead $\varphi(x)=\infty$, then
$E\varphi([e])<\infty$ and Equation~\eqref{ext:eq:total-error} forces
$\varphi(u)=\infty$. No subtraction of infinite values is used.
Since $\Fin(A)$ is nonempty, these conclusions also prove $u\ne0$.

\medskip\noindent
\textit{Apply comparison with each target.}
Zero functionals vanish at both $x$ and $u$, and the
everywhere-infinite functionals are infinite at both.
For functionals in $\Fin(A)$, the preceding estimates give
\[
 \varphi(Nu)\le\varphi(x),\qquad
 \varphi(x)\le\varphi((N+1)u).
\]
The first inequality is strict whenever $0<\varphi(x)<\infty$.
For the second, if $0<\varphi((N+1)u)<\infty$, then the functional
belongs to $\Fin(A)$ and the infinite-value case excludes
$\varphi(x)=\infty$. Equation~\eqref{ext:eq:strict-interval} makes
that inequality strict as well. Applying
Equation~\eqref{ext:eq:comparison} twice proves
$Nu\le x\le(N+1)u$.
\end{proof}

\begin{corollary}[Condition (C) from comparison]
\label{ext:cor:condition-C}
Let $A$ be a simple exact non-elementary $C^*$-algebra with strict
comparison in the sense of Equation~\eqref{ext:eq:comparison}.
Then $\Cu(A)$ satisfies condition \textup{(C)}.
\end{corollary}

\begin{proof}
Theorem~\ref{ext:thm:divisibility} gives the full almost-divisibility
sandwich, and Lemma~\ref{lem:comparison-convention} gives almost
unperforation.
\end{proof}

\begin{remark}
The same sequence $(M_j)$ works for every functional. The values
$\varphi([e])$ need not be uniformly bounded over any trace family:
their pointwise finiteness is enough to handle $\varphi(x)=\infty$,
and their bound by $\varphi(x)$ handles finite values.
This is why no trace normalization is needed.
\end{remark}

\subsection{A unital uniform-trace illustration}

The simple route into (C) is complete. We record a separate consequence
of trace transfer and integer rounding when traces can be normalized.
It illustrates the approximation mechanism without an exactness,
nuclearity, or comparison hypothesis, and is not needed in the
absorption construction.
The spectral discretization in \cite[Section~5]{DadarlatToms2010} provides
related rank-approximation methods; here small support permits integer
spectral cuts with one support-rank of error.

\begin{proposition}[Uniform approximation from below]
\label{prop:unital-approx}
Let $A$ be a unital simple infinite-dimensional $C^*$-algebra with
nonempty tracial state space $T(A)$. For every $p\in A_+$ and
$\eta>0$, there are a finite integer $L$ and $w\in M_L(A)_+$
such that
\[
 0\le\tau(p)-d_\tau(w)\le\eta
 \qquad(\tau\in T(A)),
\]
where matrix traces are unnormalized. No exactness, nuclearity, or
comparison assumption is required.
\end{proposition}

\begin{proof}
Choose an integer $N$ with $1/N<\eta$. There are $N$ mutually
orthogonal nonzero positive elements in $A$: repeatedly split a
positive element at two distinct positive spectral values. If a
nonzero hereditary corner had no such element, it would contain a
minimal projection, and simplicity would identify $A$ with compact
operators; unitality would then make $A$ finite dimensional.
For any two nonzero positives $e,f$, simplicity supplies
$0\ne v\in eAf$. The class $[v^*v]=[vv^*]$ is nonzero and
below both $[e]$ and $[f]$. Iterating over the orthogonal family
gives $0\ne h\in A_+$ with
\[
 d_\tau(h)\le 1/N\qquad(\tau\in T(A)).
\]

Fullness gives finitely many $s_i\in A$ with
$\sum_i s_i^*hs_i=1$. To see the exact equality, choose an
invertible finite sum $S=\sum_i r_iht_i$ approximating $1$.
For the row $R=(r_ih^{1/2})_i$ and column
$D=(h^{1/2}t_i)_i$, the identity $S=RD$ gives
$S^*S\le\|R\|^2D^*D$, so $D^*D$ is invertible.
Set $s_i=t_i(D^*D)^{-1/2}$.
The element
\[
 c=\sum_i h^{1/2}s_i p s_i^*h^{1/2}\in\her_A(h)_+
\]
therefore satisfies $\tau(c)=\tau(p)$ and
$d_\tau(c)\le1/N$ for every $\tau\in T(A)$.
Choose an integer $L>\|c\|$ and let
$w=\bigoplus_{j=1}^L(c-j)_+$. Lemma~\ref{ext:lem:rounding},
applied in the stabilization, gives
\[
 d_\tau(w)\le\tau(c)\le d_\tau(w)+d_\tau(c)
                       \le d_\tau(w)+1/N.
\]
This proves the required simultaneous inequalities.
\end{proof}
